\chapter{In Vitro Fertilization}

\section*{Overview}
In vitro fertilization (IVF) is an assisted reproductive technology in which eggs are retrieved from the ovaries, fertilized with sperm in the laboratory, and resulting embryos transferred to the uterus. It is the most effective form of assisted reproduction for many causes of infertility.

\section*{Alternative Names}
IVF, assisted reproductive technology, embryo transfer

\section*{Principle}
Controlled ovarian hyperstimulation with gonadotropins matures multiple follicles, which are aspirated under ultrasound guidance. Retrieved eggs are fertilized with sperm in the laboratory, and fertilized embryos are cultured for several days. One or more embryos are transferred into the uterine cavity, and any viable embryos not transferred may be frozen for later use.

\section*{History}
The first live birth from IVF, Louise Brown, occurred in 1978 in the United Kingdom. Subsequent advances included intracytoplasmic sperm injection (ICSI) in 1992 and the routine use of preimplantation genetic testing and embryo cryopreservation.

\section*{Why It Is Done}
IVF is ordered for tubal factor infertility, severe male factor infertility, advanced maternal age, endometriosis, and unexplained infertility that has failed other treatments. It is also used for preimplantation genetic testing and for fertility preservation with planned egg or embryo freezing.

\section*{Is This a Primary or Secondary Test or Procedure}
It is a primary treatment for many severe causes of infertility. It is secondary when used after failed intrauterine insemination or other fertility treatments.

\section*{How It Is Performed}
After ovulation suppression and 8-14 days of gonadotropin injections monitored by ultrasound and estradiol levels, eggs are retrieved transvaginally under sedation. Eggs are fertilized with sperm (conventional or ICSI), and embryos are cultured for 3-6 days. One or more embryos are transferred through a thin catheter, with remaining embryos often frozen.

\section*{Preparation}
Patients take daily gonadotropin injections for about 2 weeks with frequent ultrasound and blood monitoring. A full bladder is required before egg retrieval, and fasting is needed if sedation is used. Tobacco, alcohol, and some medications are discontinued.

\section*{Contraindications and Precautions}
Contraindications include active pelvic infection, severe uterine abnormality, and medical conditions making pregnancy unsafe. Precautions include screening for HIV and hepatitis, counseling on the risk of multiple gestation, and monitoring for ovarian hyperstimulation syndrome.

\section*{Risks}
Risks include ovarian hyperstimulation syndrome, multiple pregnancy, bleeding or infection from egg retrieval, and a small increase in preterm birth. The emotional and financial burden of treatment is significant.

\section*{Aftercare and Recovery}
The patient rests briefly after embryo transfer and receives progesterone support for the luteal phase. A pregnancy test is performed about 10-14 days after transfer.

\section*{Outcome and Follow-up}
A positive hCG test about 2 weeks after transfer indicates implantation, confirmed later by ultrasound. A negative result ends the cycle, and the frozen embryo bank or a new stimulation cycle may be discussed with the fertility specialist.

\section*{Expected Results}
Successful retrieval of mature eggs, normal fertilization, development of good-quality embryos, and a viable intrauterine pregnancy on follow-up ultrasound.

\section*{Follow-up}
Serial beta-hCG levels and an early ultrasound confirm a viable pregnancy, followed by routine obstetric care. Failed cycles may be followed by frozen embryo transfer or a subsequent stimulation cycle.

\section*{References}
\begin{enumerate}
  \item MedlinePlus. In Vitro Fertilization (IVF). U.S. National Library of Medicine; 2025.
  \item Current Medical Diagnosis and Treatment. McGraw-Hill; 2025.
\end{enumerate}

\clearpage
